\section{Transformasi Grammar Subset-C}
\label{sec:ch09-subsetc}

\subsection{Dari Grammar Awal ke Grammar LL(1)}

Grammar Subset-C pada Bab VIII masih memiliki dua masalah untuk
parsing top-down:
\begin{enumerate}[label=\arabic*.,leftmargin=*]
  \item \code{primary} memiliki prefiks \code{ID} yang sama dengan
        pemanggilan fungsi;
  \item \code{stmt} memiliki prefiks \code{ID} yang sama antara
        assignment dan expression statement.
\end{enumerate}

Transformasi yang dilakukan:
\begin{itemize}[leftmargin=*]
  \item memfaktorkan \code{primary} menjadi
        \code{ID [ "(" [args] ")" ]};
  \item memfaktorkan \code{stmt} menjadi
        \code{ID assign\_or\_call ";"};
  \item memisahkan assignment dari ekspresi, sehingga
        \code{expr} tidak memuat \code{ID "=" ...};
  \item mempertahankan tingkat precedence ekspresi.
\end{itemize}

\subsection{Grammar Hasil}

Grammar hasil disimpan pada \file{spec/grammar-ll1.ebnf}. Potongannya
ditunjukkan pada \lstref{lst:ch09-ll1}.

\includecode[
  language=EBNF,
  firstline=1,
  lastline=30,
  firstnumber=1,
  caption={Grammar Subset-C yang diarahkan ke LL(1)},
  label={lst:ch09-ll1}
]{../../code/subsetc/spec/grammar-ll1.ebnf}

Aturan \code{assign\_or\_call} menangani dua bentuk setelah
\code{ID}: assignment \code{= expr} atau pemanggilan
\code{( [args] )}. Dengan bentuk ini, parser hanya perlu membaca
\code{ID} sekali.

\subsection{Analisis Dampak}

\begin{table}[H]
\centering
\caption{Perubahan utama pada grammar Subset-C}
\label{tab:ch09-perubahan}
\begin{tabularx}{\textwidth}{@{}l X l@{}}
\toprule
\textbf{Aturan} & \textbf{Perubahan} & \textbf{Alasan} \\
\midrule
\code{primary} & \code{ID [ "(" [args] ")" ]} & menghilangkan prefiks
\code{ID} \\
\code{stmt} & \code{ID assign\_or\_call ";"} & memisahkan assignment
dan call \\
\code{expr} & tanpa assignment & menghindari konflik FIRST \\
\code{ifstmt} & else opsional & dangling else diselesaikan parser \\
\bottomrule
\end{tabularx}
\end{table}

\subsection{Pemeriksaan Kesetaraan Bahasa}

Setelah transformasi, penting untuk memastikan bahasa yang
dibangkitkan tidak berubah. Langkah pemeriksaan:
\begin{enumerate}[label=\arabic*.,leftmargin=*]
  \item ambil beberapa string dari grammar lama;
  \item turunkan dengan grammar baru;
  \item bandingkan pohon atau AST yang dihasilkan;
  \item uji string batas: assignment, call, if tanpa else, if dengan
        else, dan ekspresi campuran.
\end{enumerate}

Jika ada string yang dahulu valid tetapi sekarang tidak, transformasi
terlalu agresif. Jika ada string baru yang valid, transformasi
memperluas bahasa secara tidak sengaja. Keduanya harus diperbaiki
sebelum parser ditulis.

\begin{infobox}
Grammar LL(1) bukan tujuan akhir. Tujuannya adalah parser yang benar,
mudah dipelihara, dan memberikan pesan kesalahan yang baik. Jika
grammar LL(1) membuat pesan kesalahan memburuk, pertimbangkan parser
dengan lookahead lebih panjang atau pendekatan bottom-up.
\end{infobox}
